\section{Data}
\label{sec:data}

\subsection{Source and licence}
We use \emph{DataCo Smart Supply Chain for Big Data Analysis}, Version 5, by Constante, Silva
and Pereira\cite{dataco2019} (Route A, approved dataset no.~1). It is published on Mendeley Data
under a CC BY 4.0 licence and cited as the publisher requires. We use only the main table,
\tcode{DataCoSupplyChainDataset.csv}, and its column description file. The clickstream log in the
same archive is not used.

\subsection{Structure and size}
The file holds \textbf{180{,}519 order lines $\times$ 53 columns}: 65{,}752 orders (2.7 lines
each on average) from 20{,}652 customers, covering 118 products, placed between 2015-01-01 and
2018-01-31. Every line of an order carries the same label and the same order-level attributes
(the pipeline checks this), so we model at the \emph{order} level, where the decision is taken. After removing 7{,}754 cancelled
lines (\Cref{sec:prep}), \textbf{62{,}897 orders} remain. \Cref{app:dictionary} gives the data
dictionary for all 28 features used.

\subsection{Quality assessment}
\begin{itemize}
  \item \textbf{Missing values} are rare, and only in columns we do not use:
        \tcode{Product Description} is 100\% empty, \tcode{Order Zipcode} is 86.2\% empty,
        \tcode{Customer Lname} is missing in 8 rows and \tcode{Customer Zipcode} in 3. No feature
        we use has a missing value.
  \item \textbf{Outliers} were counted with the $1.5\times$IQR rule. We kept them, because
        expensive orders and heavy losses are real business events. Flagged values:
        line \tcode{Sales} 0.3\%, product price 1.1\%, and \tcode{Benefit per order} 10.5\%
        (range $-\usd{4{,}275}$ to $+\usd{912}$).
  \item \textbf{Class balance and concentration.} 57.3\% of orders are late, so predicting
        ``late'' for every order already scores 57\% accuracy. By mode (\Cref{fig:mode}): First
        Class 100\%, Second 80\%, Same Day 48\%, Standard 40\%.
  \item \textbf{Drift.} From 2017-Q4, orders shrink from about 3.0 lines to 1.0 line and the
        active catalogue falls from about 55 to 10 products. The median order value drops from
        about \usd{600} to \usd{85} (2018-Q1). Order counts and the late share stay stable
        (\Cref{fig:monthly}).
\end{itemize}

\begin{figure}[!htbp]
  \centering
  \includegraphics[width=0.85\linewidth]{01_monthly_volume_late_rate}
  \caption{Orders per month (bars, left axis) and share delivered late (line, right axis), non-cancelled orders. The dashed line marks the train/test split.}
  \label{fig:monthly}
\end{figure}

\begin{figure}[!htbp]
  \centering
  \begin{subfigure}[b]{0.47\linewidth}
    \includegraphics[width=\linewidth]{02_late_rate_by_mode}
    \caption{Late share by shipping mode (all periods).}
  \end{subfigure}\hfill
  \begin{subfigure}[b]{0.51\linewidth}
    \includegraphics[width=\linewidth]{02b_same_day_late_by_hour}
    \caption{Same Day late share by order hour (training period).}
  \end{subfigure}
  \caption{The label is concentrated by shipping mode. Within Same Day it depends only on the hour of ordering.}
  \label{fig:mode}
\end{figure}

\subsection{Is the outcome operational? A realism check}
\label{subsec:realism}
Because the label is so concentrated, we tested whether transit time depends on anything real.
The results (\Cref{tab:realism}) show it does not. Transit days are a deterministic function of
fields with no operational meaning, and order timestamps are spaced in multiples of 21
minutes. The delivery timings in this dataset were almost certainly generated by an algorithm.
The dataset is on the approved list, so we keep it, but we treat this finding as a central
limitation (\Cref{sec:limitations}). It is also the reason we compare every model against the
lookup table rather than only against a weak rule.

\begin{table}[!htbp]
  \centering
  \caption{Realism checks on the 62{,}897 orders that were not cancelled.}
  \label{tab:realism}
  \footnotesize
  \begin{tabular}{l r}
    \toprule
    \textbf{Check} & \textbf{Share of orders matching} \\
    \midrule
    Standard and Second Class: real days $= (\text{Order Id} - 1) \bmod 5 + 2$ & 100.0\% \\
    First Class: real days $= 2$                                          & 100.0\% \\
    Same Day: late $\iff$ order placed at or after 12:00                   & 100.0\% \\
    Gap between consecutive orders is a multiple of 21 minutes             & 97.5\% \\
    \bottomrule
  \end{tabular}
\end{table}
